% --- Archon physics formalization source begin ---
% archon:physics
% archon:covers PhyXMiniProblems/problem_phyx_mini_0588.lean
% archon:source-report reports/phyx_mini/problem_phyx_mini_0588.source.json

\chapter{Physics problem phyx\_mini\_0588}
\label{ch:PhyXMiniProblems_problem_phyx_mini_0588}

\paragraph{Problem source.}
\#\# Physical scenario

A light detector has an absorbing area of \textbackslash{}( 2.00\textbackslash{}times10\textasciicircum{}\{-6\}\textbackslash{},\textbackslash{}text\{m\}\textasciicircum{}2 \textbackslash{}) and absorbs 50\textbackslash{}\% of the incident light, which is at wavelength \textbackslash{}( 600\textbackslash{},\textbackslash{}text\{nm\} \textbackslash{}). The detector faces an isotropic source, \textbackslash{}( 12.0\textbackslash{},\textbackslash{}text\{m\} \textbackslash{}) from the source. The energy \textbackslash{}( E \textbackslash{}) emitted by the source versus time \textbackslash{}( t \textbackslash{}) is given in figure \textbackslash{}( (E\_s = 7.2\textbackslash{},\textbackslash{}text\{nJ\},\textbackslash{},t\_s = 2.0\textbackslash{},\textbackslash{}text\{s\}) \textbackslash{}).

\#\# Question

At what rate are photons absorbed by the detector?

\#\# Answer choices

- A: A: 12.0photons/s

- B: B: 3.0photons/s

- C: C: 6.0photons/s

- D: D: 9.0photons/s

\#\# Figure caption (auxiliary; use the image as primary evidence)

The image is a graph with the x-axis labeled as \textbackslash{}( t \textbackslash{}) (s) and the y-axis labeled as \textbackslash{}( E \textbackslash{}) (nJ). There is a diagonal line in pink that starts from the origin (0, 0) and rises linearly up to the point \textbackslash{}( (t\_s, E\_s) \textbackslash{}). At this point, two dashed lines extend horizontally and vertically: one from \textbackslash{}( E\_s \textbackslash{}) to the y-axis and another from \textbackslash{}( t\_s \textbackslash{}) to the x-axis. This indicates a relationship where energy \textbackslash{}( E \textbackslash{}) increases linearly over time \textbackslash{}( t \textbackslash{}) until it reaches the specified energy level \textbackslash{}( E\_s \textbackslash{}) at time \textbackslash{}( t\_s \textbackslash{}).

\#\# Dataset metadata

Quantum Phenomena; Multi-Formula Reasoning, Physical Model Grounding Reasoning

\paragraph{Recorded answer/context.}
C: C: 6.0photons/s

\paragraph{Figure/image path.}
/root/proposal\_for\_physic/science-mango/phyx\_mini\_run/phyx\_data/test\_image/588.png

\paragraph{Formalization target.}
create a compiling Lean file with sorry bodies at `PhyXMiniProblems/problem\_phyx\_mini\_0588.lean`. The Lean declarations must preserve the physical quantities, dimensions or dimensional roles, figure labels, governing-law hypotheses, and final relation expressed by this problem.
Use Mathlib/Physlib names found through LeanExplore where available. If a physics API is missing, introduce faithful local abstractions rather than scalar placeholder aliases.

\begin{theorem}[Physics formalization target]
\label{thm:physics:phyx_mini_0588:target}
\lean{PhyXMiniProblems.ProblemPhyXMini0588.problem_phyx_mini_0588}
\uses{def:physics:phyx-mini-0588:phyxminiproblems-problemphyxmini0588-photondetectorsetup, def:physics:phyx-mini-0588:phyxminiproblems-problemphyxmini0588-matchesphotondetectorscenario, def:physics:phyx-mini-0588:phyxminiproblems-problemphyxmini0588-matchesproblemreadouts, def:physics:phyx-mini-0588:phyxminiproblems-problemphyxmini0588-matchessourceenergyfigure, def:physics:phyx-mini-0588:phyxminiproblems-problemphyxmini0588-matchescalibratedconstants, def:physics:phyx-mini-0588:phyxminiproblems-problemphyxmini0588-photonabsorptionlaws, def:physics:phyx-mini-0588:phyxminiproblems-problemphyxmini0588-matchesdisplayedphotonrate, def:physics:phyx-mini-0588:phyxminiproblems-problemphyxmini0588-isuniquematchingphotonrate}
The assigned autoformalize agent should translate this physics problem into Lean declarations in the covered file, with theorem and lemma proof bodies written as `by sorry`.
\end{theorem}
\begin{proof}
This is an autoformalization task, not a proof task. Produce statements that can later be proved by the physics prover without weakening the physical model.
\end{proof}

% --- Archon declaration topology begin ---
\section{Declaration topology}
\begin{definition}[power Dimension]
  \label{def:physics:phyx-mini-0588:phyxminiproblems-problemphyxmini0588-powerdimension}
  \lean{PhyXMiniProblems.ProblemPhyXMini0588.powerDimension}
  Power has physical dimension energy divided by time
\end{definition}
\begin{proof}
  This is the declared model definition of power Dimension.
\end{proof}
\begin{definition}[action Dimension]
  \label{def:physics:phyx-mini-0588:phyxminiproblems-problemphyxmini0588-actiondimension}
  \lean{PhyXMiniProblems.ProblemPhyXMini0588.actionDimension}
  Action, the dimensional role of ordinary Planck's constant, is energy-time
\end{definition}
\begin{proof}
  This is the declared model definition of action Dimension.
\end{proof}
\begin{definition}[Length Quantity]
  \label{def:physics:phyx-mini-0588:phyxminiproblems-problemphyxmini0588-lengthquantity}
  \lean{PhyXMiniProblems.ProblemPhyXMini0588.LengthQuantity}
  A nonnegative, unit-independent physical length
\end{definition}
\begin{proof}
  This is the declared model definition of Length Quantity.
\end{proof}
\begin{definition}[Time Quantity]
  \label{def:physics:phyx-mini-0588:phyxminiproblems-problemphyxmini0588-timequantity}
  \lean{PhyXMiniProblems.ProblemPhyXMini0588.TimeQuantity}
  A nonnegative, unit-independent physical time
\end{definition}
\begin{proof}
  This is the declared model definition of Time Quantity.
\end{proof}
\begin{definition}[Power Quantity]
  \label{def:physics:phyx-mini-0588:phyxminiproblems-problemphyxmini0588-powerquantity}
  \lean{PhyXMiniProblems.ProblemPhyXMini0588.PowerQuantity}
  \uses{def:physics:phyx-mini-0588:phyxminiproblems-problemphyxmini0588-powerdimension}
  A nonnegative, unit-independent physical power
\end{definition}
\begin{proof}
  This is the declared model definition of Power Quantity.
\end{proof}
\begin{definition}[Planck Constant Quantity]
  \label{def:physics:phyx-mini-0588:phyxminiproblems-problemphyxmini0588-planckconstantquantity}
  \lean{PhyXMiniProblems.ProblemPhyXMini0588.PlanckConstantQuantity}
  \uses{def:physics:phyx-mini-0588:phyxminiproblems-problemphyxmini0588-actiondimension}
  A nonnegative, unit-independent value of ordinary Planck's constant
\end{definition}
\begin{proof}
  This is the declared model definition of Planck Constant Quantity.
\end{proof}
\begin{definition}[Photon Rate Quantity]
  \label{def:physics:phyx-mini-0588:phyxminiproblems-problemphyxmini0588-photonratequantity}
  \lean{PhyXMiniProblems.ProblemPhyXMini0588.PhotonRateQuantity}
  A nonnegative photon-count rate, with dimension inverse time
\end{definition}
\begin{proof}
  This is the declared model definition of Photon Rate Quantity.
\end{proof}
\begin{definition}[nonnegative SIReadout]
  \label{def:physics:phyx-mini-0588:phyxminiproblems-problemphyxmini0588-nonnegativesireadout}
  \lean{PhyXMiniProblems.ProblemPhyXMini0588.nonnegativeSIReadout}
  Coherent-SI readout of a nonnegative dimensionful quantity
\end{definition}
\begin{proof}
  This is the declared model definition of nonnegative SIReadout.
\end{proof}
\begin{definition}[area In Square Meters]
  \label{def:physics:phyx-mini-0588:phyxminiproblems-problemphyxmini0588-areainsquaremeters}
  \lean{PhyXMiniProblems.ProblemPhyXMini0588.areaInSquareMeters}
  \uses{def:physics:phyx-mini-0588:phyxminiproblems-problemphyxmini0588-nonnegativesireadout}
  Area readout in square metres
\end{definition}
\begin{proof}
  This is the declared model definition of area In Square Meters.
\end{proof}
\begin{definition}[length In Meters]
  \label{def:physics:phyx-mini-0588:phyxminiproblems-problemphyxmini0588-lengthinmeters}
  \lean{PhyXMiniProblems.ProblemPhyXMini0588.lengthInMeters}
  \uses{def:physics:phyx-mini-0588:phyxminiproblems-problemphyxmini0588-lengthquantity, def:physics:phyx-mini-0588:phyxminiproblems-problemphyxmini0588-nonnegativesireadout}
  Length readout in metres
\end{definition}
\begin{proof}
  This is the declared model definition of length In Meters.
\end{proof}
\begin{definition}[length In Nanometers]
  \label{def:physics:phyx-mini-0588:phyxminiproblems-problemphyxmini0588-lengthinnanometers}
  \lean{PhyXMiniProblems.ProblemPhyXMini0588.lengthInNanometers}
  \uses{def:physics:phyx-mini-0588:phyxminiproblems-problemphyxmini0588-lengthquantity, def:physics:phyx-mini-0588:phyxminiproblems-problemphyxmini0588-lengthinmeters}
  Length readout in nanometres
\end{definition}
\begin{proof}
  This is the declared model definition of length In Nanometers.
\end{proof}
\begin{definition}[time In Seconds]
  \label{def:physics:phyx-mini-0588:phyxminiproblems-problemphyxmini0588-timeinseconds}
  \lean{PhyXMiniProblems.ProblemPhyXMini0588.timeInSeconds}
  \uses{def:physics:phyx-mini-0588:phyxminiproblems-problemphyxmini0588-timequantity, def:physics:phyx-mini-0588:phyxminiproblems-problemphyxmini0588-nonnegativesireadout}
  Time readout in seconds
\end{definition}
\begin{proof}
  This is the declared model definition of time In Seconds.
\end{proof}
\begin{definition}[energy In Joules]
  \label{def:physics:phyx-mini-0588:phyxminiproblems-problemphyxmini0588-energyinjoules}
  \lean{PhyXMiniProblems.ProblemPhyXMini0588.energyInJoules}
  Energy readout in joules
\end{definition}
\begin{proof}
  This is the declared model definition of energy In Joules.
\end{proof}
\begin{definition}[energy In Nanojoules]
  \label{def:physics:phyx-mini-0588:phyxminiproblems-problemphyxmini0588-energyinnanojoules}
  \lean{PhyXMiniProblems.ProblemPhyXMini0588.energyInNanojoules}
  \uses{def:physics:phyx-mini-0588:phyxminiproblems-problemphyxmini0588-energyinjoules}
  Energy readout in nanojoules
\end{definition}
\begin{proof}
  This is the declared model definition of energy In Nanojoules.
\end{proof}
\begin{definition}[power In Watts]
  \label{def:physics:phyx-mini-0588:phyxminiproblems-problemphyxmini0588-powerinwatts}
  \lean{PhyXMiniProblems.ProblemPhyXMini0588.powerInWatts}
  \uses{def:physics:phyx-mini-0588:phyxminiproblems-problemphyxmini0588-powerquantity, def:physics:phyx-mini-0588:phyxminiproblems-problemphyxmini0588-nonnegativesireadout}
  Power readout in watts
\end{definition}
\begin{proof}
  This is the declared model definition of power In Watts.
\end{proof}
\begin{definition}[speed In Meters Per Second]
  \label{def:physics:phyx-mini-0588:phyxminiproblems-problemphyxmini0588-speedinmeterspersecond}
  \lean{PhyXMiniProblems.ProblemPhyXMini0588.speedInMetersPerSecond}
  \uses{def:physics:phyx-mini-0588:phyxminiproblems-problemphyxmini0588-nonnegativesireadout}
  Speed readout in metres per second
\end{definition}
\begin{proof}
  This is the declared model definition of speed In Meters Per Second.
\end{proof}
\begin{definition}[planck Constant In Joule Seconds]
  \label{def:physics:phyx-mini-0588:phyxminiproblems-problemphyxmini0588-planckconstantinjouleseconds}
  \lean{PhyXMiniProblems.ProblemPhyXMini0588.planckConstantInJouleSeconds}
  \uses{def:physics:phyx-mini-0588:phyxminiproblems-problemphyxmini0588-planckconstantquantity, def:physics:phyx-mini-0588:phyxminiproblems-problemphyxmini0588-nonnegativesireadout}
  Ordinary Planck-constant readout in joule-seconds
\end{definition}
\begin{proof}
  This is the declared model definition of planck Constant In Joule Seconds.
\end{proof}
\begin{definition}[photon Rate In Per Second]
  \label{def:physics:phyx-mini-0588:phyxminiproblems-problemphyxmini0588-photonrateinpersecond}
  \lean{PhyXMiniProblems.ProblemPhyXMini0588.photonRateInPerSecond}
  \uses{def:physics:phyx-mini-0588:phyxminiproblems-problemphyxmini0588-photonratequantity, def:physics:phyx-mini-0588:phyxminiproblems-problemphyxmini0588-nonnegativesireadout}
  Photon-count rate readout in photons per second
\end{definition}
\begin{proof}
  This is the declared model definition of photon Rate In Per Second.
\end{proof}
\begin{definition}[Emission Pattern]
  \label{def:physics:phyx-mini-0588:phyxminiproblems-problemphyxmini0588-emissionpattern}
  \lean{PhyXMiniProblems.ProblemPhyXMini0588.EmissionPattern}
  Angular emission model of the light source
\end{definition}
\begin{proof}
  This is the declared model definition of Emission Pattern.
\end{proof}
\begin{definition}[Light Spectrum]
  \label{def:physics:phyx-mini-0588:phyxminiproblems-problemphyxmini0588-lightspectrum}
  \lean{PhyXMiniProblems.ProblemPhyXMini0588.LightSpectrum}
  Spectral model of the incident light
\end{definition}
\begin{proof}
  This is the declared model definition of Light Spectrum.
\end{proof}
\begin{definition}[Detector Orientation]
  \label{def:physics:phyx-mini-0588:phyxminiproblems-problemphyxmini0588-detectororientation}
  \lean{PhyXMiniProblems.ProblemPhyXMini0588.DetectorOrientation}
  Orientation of the detector's absorbing face
\end{definition}
\begin{proof}
  This is the declared model definition of Detector Orientation.
\end{proof}
\begin{definition}[Figure Axis]
  \label{def:physics:phyx-mini-0588:phyxminiproblems-problemphyxmini0588-figureaxis}
  \lean{PhyXMiniProblems.ProblemPhyXMini0588.FigureAxis}
  The two axes visible in the source-energy graph
\end{definition}
\begin{proof}
  This is the declared model definition of Figure Axis.
\end{proof}
\begin{definition}[Figure Axis Quantity]
  \label{def:physics:phyx-mini-0588:phyxminiproblems-problemphyxmini0588-figureaxisquantity}
  \lean{PhyXMiniProblems.ProblemPhyXMini0588.FigureAxisQuantity}
  Physical quantity assigned to an axis of the graph
\end{definition}
\begin{proof}
  This is the declared model definition of Figure Axis Quantity.
\end{proof}
\begin{definition}[Figure Axis Unit]
  \label{def:physics:phyx-mini-0588:phyxminiproblems-problemphyxmini0588-figureaxisunit}
  \lean{PhyXMiniProblems.ProblemPhyXMini0588.FigureAxisUnit}
  Unit printed beside an axis of the graph
\end{definition}
\begin{proof}
  This is the declared model definition of Figure Axis Unit.
\end{proof}
\begin{definition}[Figure Trace Color]
  \label{def:physics:phyx-mini-0588:phyxminiproblems-problemphyxmini0588-figuretracecolor}
  \lean{PhyXMiniProblems.ProblemPhyXMini0588.FigureTraceColor}
  Color of the rising trace in the primary raster
\end{definition}
\begin{proof}
  This is the declared model definition of Figure Trace Color.
\end{proof}
\begin{definition}[Source Energy Figure]
  \label{def:physics:phyx-mini-0588:phyxminiproblems-problemphyxmini0588-sourceenergyfigure}
  \lean{PhyXMiniProblems.ProblemPhyXMini0588.SourceEnergyFigure}
  \uses{def:physics:phyx-mini-0588:phyxminiproblems-problemphyxmini0588-timequantity, def:physics:phyx-mini-0588:phyxminiproblems-problemphyxmini0588-figureaxis, def:physics:phyx-mini-0588:phyxminiproblems-problemphyxmini0588-figureaxisquantity, def:physics:phyx-mini-0588:phyxminiproblems-problemphyxmini0588-figureaxisunit, def:physics:phyx-mini-0588:phyxminiproblems-problemphyxmini0588-figuretracecolor}
  The primary figure. Its labeled point carries physical time and energy quantities; the booleans retain the origin, straight-line, and dashed-guide features visible in the raster
\end{definition}
\begin{proof}
  This is the declared model definition of Source Energy Figure.
\end{proof}
\begin{definition}[Photon Detector Setup]
  \label{def:physics:phyx-mini-0588:phyxminiproblems-problemphyxmini0588-photondetectorsetup}
  \lean{PhyXMiniProblems.ProblemPhyXMini0588.PhotonDetectorSetup}
  \uses{def:physics:phyx-mini-0588:phyxminiproblems-problemphyxmini0588-lengthquantity, def:physics:phyx-mini-0588:phyxminiproblems-problemphyxmini0588-powerquantity, def:physics:phyx-mini-0588:phyxminiproblems-problemphyxmini0588-planckconstantquantity, def:physics:phyx-mini-0588:phyxminiproblems-problemphyxmini0588-photonratequantity, def:physics:phyx-mini-0588:phyxminiproblems-problemphyxmini0588-emissionpattern, def:physics:phyx-mini-0588:phyxminiproblems-problemphyxmini0588-lightspectrum, def:physics:phyx-mini-0588:phyxminiproblems-problemphyxmini0588-detectororientation, def:physics:phyx-mini-0588:phyxminiproblems-problemphyxmini0588-sourceenergyfigure}
  Independent physical quantities in the experiment. The four rate/energy fields are observables related by the governing laws below, not definitions of the desired numerical answer
\end{definition}
\begin{proof}
  This is the declared model definition of Photon Detector Setup.
\end{proof}
\begin{definition}[Matches Photon Detector Scenario]
  \label{def:physics:phyx-mini-0588:phyxminiproblems-problemphyxmini0588-matchesphotondetectorscenario}
  \lean{PhyXMiniProblems.ProblemPhyXMini0588.MatchesPhotonDetectorScenario}
  \uses{def:physics:phyx-mini-0588:phyxminiproblems-problemphyxmini0588-photondetectorsetup}
  Qualitative roles stated in the prose description
\end{definition}
\begin{proof}
  This is the declared model definition of Matches Photon Detector Scenario.
\end{proof}
\begin{definition}[Matches Problem Readouts]
  \label{def:physics:phyx-mini-0588:phyxminiproblems-problemphyxmini0588-matchesproblemreadouts}
  \lean{PhyXMiniProblems.ProblemPhyXMini0588.MatchesProblemReadouts}
  \uses{def:physics:phyx-mini-0588:phyxminiproblems-problemphyxmini0588-areainsquaremeters, def:physics:phyx-mini-0588:phyxminiproblems-problemphyxmini0588-lengthinmeters, def:physics:phyx-mini-0588:phyxminiproblems-problemphyxmini0588-lengthinnanometers, def:physics:phyx-mini-0588:phyxminiproblems-problemphyxmini0588-photondetectorsetup}
  Scalar readouts given in the problem prose. These premises state detector geometry, absorptivity, wavelength, and distance; they contain no photon-rate conclusion
\end{definition}
\begin{proof}
  This is the declared model definition of Matches Problem Readouts.
\end{proof}
\begin{definition}[Matches Source Energy Figure]
  \label{def:physics:phyx-mini-0588:phyxminiproblems-problemphyxmini0588-matchessourceenergyfigure}
  \lean{PhyXMiniProblems.ProblemPhyXMini0588.MatchesSourceEnergyFigure}
  \uses{def:physics:phyx-mini-0588:phyxminiproblems-problemphyxmini0588-timeinseconds, def:physics:phyx-mini-0588:phyxminiproblems-problemphyxmini0588-energyinjoules, def:physics:phyx-mini-0588:phyxminiproblems-problemphyxmini0588-energyinnanojoules, def:physics:phyx-mini-0588:phyxminiproblems-problemphyxmini0588-photondetectorsetup}
  Literal graph labels and qualitative evidence from image 588. The final field expresses the straight trace as proportional growth between the origin and the labeled point, without assuming its slope or the requested photon rate
\end{definition}
\begin{proof}
  This is the declared model definition of Matches Source Energy Figure.
\end{proof}
\begin{definition}[Matches Calibrated Constants]
  \label{def:physics:phyx-mini-0588:phyxminiproblems-problemphyxmini0588-matchescalibratedconstants}
  \lean{PhyXMiniProblems.ProblemPhyXMini0588.MatchesCalibratedConstants}
  \uses{def:physics:phyx-mini-0588:phyxminiproblems-problemphyxmini0588-speedinmeterspersecond, def:physics:phyx-mini-0588:phyxminiproblems-problemphyxmini0588-planckconstantinjouleseconds, def:physics:phyx-mini-0588:phyxminiproblems-problemphyxmini0588-photondetectorsetup}
  Exact coherent-SI readouts of the two universal constants used by the photon model. They are calibrated physical data, not consequences about this detector
\end{definition}
\begin{proof}
  This is the declared model definition of Matches Calibrated Constants.
\end{proof}
\begin{definition}[Photon Absorption Laws]
  \label{def:physics:phyx-mini-0588:phyxminiproblems-problemphyxmini0588-photonabsorptionlaws}
  \lean{PhyXMiniProblems.ProblemPhyXMini0588.PhotonAbsorptionLaws}
  \uses{def:physics:phyx-mini-0588:phyxminiproblems-problemphyxmini0588-areainsquaremeters, def:physics:phyx-mini-0588:phyxminiproblems-problemphyxmini0588-lengthinmeters, def:physics:phyx-mini-0588:phyxminiproblems-problemphyxmini0588-timeinseconds, def:physics:phyx-mini-0588:phyxminiproblems-problemphyxmini0588-energyinjoules, def:physics:phyx-mini-0588:phyxminiproblems-problemphyxmini0588-powerinwatts, def:physics:phyx-mini-0588:phyxminiproblems-problemphyxmini0588-speedinmeterspersecond, def:physics:phyx-mini-0588:phyxminiproblems-problemphyxmini0588-planckconstantinjouleseconds, def:physics:phyx-mini-0588:phyxminiproblems-problemphyxmini0588-photonrateinpersecond, def:physics:phyx-mini-0588:phyxminiproblems-problemphyxmini0588-photondetectorsetup}
  The five physical relations used in the calculation are: graph slope equals source radiant power; isotropic power crosses a sphere of area 4 π r² uniformly; absorbed power is absorptivity times incident power; Eγ λ = h c; and photon energy balance converts power to count rate. None of these fields fixes the count rate to a displayed answer value
\end{definition}
\begin{proof}
  This is the declared model definition of Photon Absorption Laws.
\end{proof}
\begin{definition}[Answer Choice]
  \label{def:physics:phyx-mini-0588:phyxminiproblems-problemphyxmini0588-answerchoice}
  \lean{PhyXMiniProblems.ProblemPhyXMini0588.AnswerChoice}
  Answer labels in the order printed in the problem
\end{definition}
\begin{proof}
  This is the declared model definition of Answer Choice.
\end{proof}
\begin{definition}[displayed Photon Rate In Per Second]
  \label{def:physics:phyx-mini-0588:phyxminiproblems-problemphyxmini0588-displayedphotonrateinpersecond}
  \lean{PhyXMiniProblems.ProblemPhyXMini0588.displayedPhotonRateInPerSecond}
  \uses{def:physics:phyx-mini-0588:phyxminiproblems-problemphyxmini0588-answerchoice}
  Displayed photon rates, in photons per second
\end{definition}
\begin{proof}
  This is the declared model definition of displayed Photon Rate In Per Second.
\end{proof}
\begin{definition}[displayed Rate Tolerance]
  \label{def:physics:phyx-mini-0588:phyxminiproblems-problemphyxmini0588-displayedratetolerance}
  \lean{PhyXMiniProblems.ProblemPhyXMini0588.displayedRateTolerance}
  Tolerance corresponding to a rate displayed to one decimal place
\end{definition}
\begin{proof}
  This is the declared model definition of displayed Rate Tolerance.
\end{proof}
\begin{definition}[Matches Displayed Photon Rate]
  \label{def:physics:phyx-mini-0588:phyxminiproblems-problemphyxmini0588-matchesdisplayedphotonrate}
  \lean{PhyXMiniProblems.ProblemPhyXMini0588.MatchesDisplayedPhotonRate}
  \uses{def:physics:phyx-mini-0588:phyxminiproblems-problemphyxmini0588-photonrateinpersecond, def:physics:phyx-mini-0588:phyxminiproblems-problemphyxmini0588-photondetectorsetup, def:physics:phyx-mini-0588:phyxminiproblems-problemphyxmini0588-answerchoice, def:physics:phyx-mini-0588:phyxminiproblems-problemphyxmini0588-displayedphotonrateinpersecond, def:physics:phyx-mini-0588:phyxminiproblems-problemphyxmini0588-displayedratetolerance}
  The computed physical rate rounds to the displayed value of a choice
\end{definition}
\begin{proof}
  This is the declared model definition of Matches Displayed Photon Rate.
\end{proof}
\begin{definition}[Is Unique Matching Photon Rate]
  \label{def:physics:phyx-mini-0588:phyxminiproblems-problemphyxmini0588-isuniquematchingphotonrate}
  \lean{PhyXMiniProblems.ProblemPhyXMini0588.IsUniqueMatchingPhotonRate}
  \uses{def:physics:phyx-mini-0588:phyxminiproblems-problemphyxmini0588-photondetectorsetup, def:physics:phyx-mini-0588:phyxminiproblems-problemphyxmini0588-answerchoice, def:physics:phyx-mini-0588:phyxminiproblems-problemphyxmini0588-matchesdisplayedphotonrate}
  Exactly one printed alternative matches the computed physical rate
\end{definition}
\begin{proof}
  This is the declared model definition of Is Unique Matching Photon Rate.
\end{proof}
% --- Archon declaration topology end ---
% --- Archon physics formalization source end ---
